\chapter{Proof-of-Concept Artifacts}
\label{appendix:pocs}

The research produced a set of proof-of-concept programs that demonstrate the leaks and their
reach. As the Technion guidelines ask that software developed for a thesis be described in an
appendix, we summarize them here. We describe their structure and attacker model. Their source is
available at \url{https://github.com/GalKaptsenel/thesis-artifacts}. Each is \emph{self-scoring}:
the victim draws a fresh random secret, which the attacker then retrieves, and afterwards they are
compared for scoring. All measurements are on the Neoverse~N3, pinned to isolated cores. The attacker runs
unprivileged at EL0. No elevated privilege, special instruction, or configuration is mandatory
beyond a vulnerable code pattern and co-location on a core.

\section{Cross-boundary PoC}
\label{app:poc-xboundary}

\reg{cross\_boundary --switch \{aut|canon\} --boundary \{same|xthread|xproc|xel|xvm\}} recovers a
random pointer-authentication signature, used as the target secret and regenerated each round,
across an isolation boundary (\autoref{sec:eval-reach}). Each round the victim signs a pointer and,
inside one Spectre-v4 store-bypass window, manipulates it according to the chosen covert channel and
uses it. The \reg{canon} switch masks a single bit, so the pointer's canonicality corresponds to the
leaked bit value. The \reg{aut} switch places the leaked bit at bit~55 and authenticates, so sign
extension leaves the pointer in the user half or the kernel half. The read-out runs over the
requested medium (\reg{--medium}: L1D, data or instruction TLB, or BTB), timed by a PMU event
counter or the low-privilege \reg{cntvct\_el0} virtual counter (\reg{--timer}).

\section{Blind stack-canary recovery}
\label{app:poc-canary}

\reg{canary\_leak\_blind -w \{v1|v4|rsb\}} recovers a process's own stack canary from a
two-dimensional-array gadget driven by Spectre-v1, Spectre-v4, or Spectre-RSB, reading the
canonicality oracle through Flush+Reload~\cite{flushreload} over a pre-mapped array of probe pages
(\autoref{sec:eval-exploit}). It combines an additive sweep for the top bits with a top-down carry
propagation for the lower bits, and recovers the full \reg{[55:0]} value.

\section{RSA-2048 key recovery through OpenSSL}
\label{app:poc-openssl}

\reg{rsa\_leak} recovers the bytes of an in-use RSA-2048 private key from OpenSSL built with its ROP
mitigations disabled and a single planted vulnerable \texttt{memcpy} (\autoref{sec:eval-exploit}). A
stack-buffer overflow drives a return-oriented chain built entirely from OpenSSL's own gadgets,
fires a Spectre-v4 window, and transmits over the BTB channel read with \reg{cntvct\_el0}.

\section{An EL3 round trip does not flush the BTB}
\label{app:poc-smc}

\reg{smc\_bench} tests a condition for crossing the secure boundary. It primes one EL0 BTB entry,
performs an EL3 round trip through a Secure Monitor Call (SMC) between prime and probe, and reports
that the entry survives, so an EL3 world switch does not by itself destroy the prediction state the
BTB channel relies on. This indicates the address-validity leak's BTB channel is feasible across
secure-world switches.
